
\documentclass{article}
\usepackage{colortbl}
\usepackage{makecell}
\usepackage{multirow}
\usepackage{supertabular}

\begin{document}

\newcounter{utterance}

\twocolumn

{ \footnotesize  \setcounter{utterance}{1}
\setlength{\tabcolsep}{0pt}
\begin{supertabular}{c@{$\;$}|p{.15\linewidth}@{}p{.15\linewidth}p{.15\linewidth}p{.15\linewidth}p{.15\linewidth}p{.15\linewidth}}

    \# & $\;$A & \multicolumn{4}{c}{Game Master} & $\;\:$B\\
    \hline 

    \theutterance \stepcounter{utterance}  

    & & \multicolumn{4}{p{0.6\linewidth}}{\cellcolor[rgb]{0.9,0.9,0.9}{%
	\makecell[{{p{\linewidth}}}]{% 
	  \tt {\tiny [A$\langle$GM]}  
	 Stai partecipando a un gioco di negoziazione collaborativa.\\ \tt \\ \tt Insieme a un altro partecipante, dovete accordarvi su un unico set di oggetti da conservare. Ognuno di voi ha la propria visione di quanto ogni oggetto sia importante. Non sai come l'altro partecipante valuta gli oggetti. Inoltre, ti viene fornito lo sforzo che ogni oggetto richiede.\\ \tt Potete accordarvi su un set solo se lo sforzo totale degli oggetti selezionati non supera un limite condiviso:\\ \tt \\ \tt LIMITE = 4026\\ \tt \\ \tt Ecco i valori di sforzo dei singoli oggetti:\\ \tt \\ \tt Sforzo oggetto = {"C17": 780, "B41": 742, "A71": 997, "C00": 665, "B35": 140, "A77": 77, "A03": 513, "A04": 383, "C54": 587, "C01": 825, "A87": 319, "B32": 448, "B86": 516, "B33": 694, "A10": 366}\\ \tt \\ \tt Ecco la tua visione personale sull'importanza di ogni oggetto:\\ \tt \\ \tt Valori di importanza dell'oggetto = {"C17": 138, "B41": 583, "A71": 868, "C00": 822, "B35": 783, "A77": 65, "A03": 262, "A04": 121, "C54": 508, "C01": 780, "A87": 461, "B32": 484, "B86": 668, "B33": 389, "A10": 808}\\ \tt \\ \tt Obiettivo:\\ \tt \\ \tt Il tuo obiettivo è negoziare un set condiviso di oggetti che ti avvantaggi il più possibile (cioè, massimizzi l'importanza totale per TE), rimanendo entro il LIMITEE. Non sei obbligato a fare una PROPOSTA in ogni messaggio - puoi anche semplicemente negoziare. Tutte le tattiche sono permesse!\\ \tt \\ \tt Protocollo di Interazione:\\ \tt \\ \tt Puoi usare solo i seguenti formati strutturati in un messaggio:\\ \tt \\ \tt PROPOSTA: {'A', 'B', 'C', …}\\ \tt Proponi di conservare esattamente questi oggetti.\\ \tt RIFIUTO: {'A', 'B', 'C', …}\\ \tt Rifiuta esplicitamente la proposta dell'avversario.\\ \tt ARGOMENTAZIONE: {'...'}\\ \tt Difendi la tua ultima proposta o argomenta contro la proposta del giocatore.\\ \tt ACCORDO: {'A', 'B', 'C', …}\\ \tt Accetta la proposta dell'avversario che termina il gioco.\\ \tt RAGIONAMENTO STRATEGICO: {'...'}\\ \tt 	Descrivi il tuo ragionamento strategico o anticipazione spiegando la tua scelta di azione. Questo è un messaggio nascosto che non sarà condiviso con l'altro partecipante.\\ \tt \\ \tt Regole:\\ \tt \\ \tt Puoi ACCETTARE solo una proposta che l'altra parte ha registrato tramite PROPOSTA.\\ \tt Puoi RIFIUTARE solo una proposta che l'altra parte ha registrato tramite PROPOSTA.\\ \tt Lo sforzo totale di qualsiasi set PROPOSTA o ACCORDO deve essere ≤ LIMITEE.\\ \tt NON rivelare i tuoi punteggi di importanza nascosti.\\ \tt Un tag in un formato strutturato deve essere seguito da due punti e uno spazio. L'argomento deve essere un set python contenente 0 o più stringhe.\\ \tt Quindi, deve essere della forma TAG: {...}\\ \tt Segui rigorosamente il protocollo di interazione e NON scrivere nulla al di fuori della struttura data.\\ \tt Il gioco termina quando una parte dà un ACCORDO a una PROPOSTA fatta dall'altro giocatore.\\ \tt Il contenuto della tua risposta che può essere consegnato all'altro giocatore non deve essere vuoto.\\ \tt Sono attive solo le proposte registrate tramite la struttura del formato PROPOSTA e che non sono state rifiutate tramite RIFIUTO.\\ \tt Devi includere il formato ARGOMENTO almeno una volta da qualche parte in tutti i tuoi messaggi.\\ \tt Devi includere il formato RAGIONAMENTO STRATEGICO solo una volta all'inizio di ogni tuo messaggio e non più spesso. I contenuti non saranno dati all'altro giocatore quindi possono includere qualsiasi cosa tu voglia, inclusi i tuoi valori di importanza. Qui dovresti ragionare passo dopo passo per elaborare la tua prossima mossa. 
	  } 
	   } 
	   } 
	 & \\ 
 

    \theutterance \stepcounter{utterance}  

    & \multicolumn{4}{p{0.6\linewidth}}{\cellcolor[rgb]{0.8,1,0.9}{%
	\makecell[{{p{\linewidth}}}]{% 
	  \tt {\tiny [A$\rangle$GM]}  
	 RAGIONAMENTO STRATEGICO: {'Sto iniziando il gioco e non ho informazioni sulla valutazione dell'altro giocatore, quindi devo iniziare proponendo un set di oggetti che massimizzino l'importanza totale per me, rimanendo entro il limite di sforzo. Inizierò con gli oggetti più importanti per me e poi aggiusterò la proposta in base alle risposte dell'altro giocatore.'}\\ \tt \\ \tt ARGOMENTAZIONE: {'Inizierò proponendo un set di oggetti che includono i più importanti per me, come A71, C00 e B41, che hanno un'importanza totale di 868 + 822 + 583 = 2273. Lo sforzo totale per questi oggetti è 997 + 665 + 742 = 2404, che è inferiore al limite di 4026. Questa proposta mi permette di massimizzare l'importanza totale per me e rimanere entro il limite di sforzo.'}\\ \tt \\ \tt PROPOSTA: {'A71', 'C00', 'B41'} 
	  } 
	   } 
	   } 
	 & & \\ 
 

    \theutterance \stepcounter{utterance}  

    & & \multicolumn{4}{p{0.6\linewidth}}{\cellcolor[rgb]{0.9,0.9,0.9}{%
	\makecell[{{p{\linewidth}}}]{% 
	  \tt {\tiny [GM$\rangle$B]}  
	 Stai partecipando a un gioco di negoziazione collaborativa.\\ \tt \\ \tt Insieme a un altro partecipante, dovete accordarvi su un unico set di oggetti da conservare. Ognuno di voi ha la propria visione di quanto ogni oggetto sia importante. Non sai come l'altro partecipante valuta gli oggetti. Inoltre, ti viene fornito lo sforzo che ogni oggetto richiede.\\ \tt Potete accordarvi su un set solo se lo sforzo totale degli oggetti selezionati non supera un limite condiviso:\\ \tt \\ \tt LIMITE = 4026\\ \tt \\ \tt Ecco i valori di sforzo dei singoli oggetti:\\ \tt \\ \tt Sforzo oggetto = {"C17": 780, "B41": 742, "A71": 997, "C00": 665, "B35": 140, "A77": 77, "A03": 513, "A04": 383, "C54": 587, "C01": 825, "A87": 319, "B32": 448, "B86": 516, "B33": 694, "A10": 366}\\ \tt \\ \tt Ecco la tua visione personale sull'importanza di ogni oggetto:\\ \tt \\ \tt Valori di importanza dell'oggetto = {"A71": 65, "C00": 121, "A10": 138, "B35": 262, "C01": 389, "B86": 461, "B41": 484, "C54": 508, "B32": 583, "A87": 668, "B33": 780, "A03": 783, "C17": 808, "A04": 822, "A77": 868}\\ \tt \\ \tt Obiettivo:\\ \tt \\ \tt Il tuo obiettivo è negoziare un set condiviso di oggetti che ti avvantaggi il più possibile (cioè, massimizzi l'importanza totale per TE), rimanendo entro il LIMITEE. Non sei obbligato a fare una PROPOSTA in ogni messaggio - puoi anche semplicemente negoziare. Tutte le tattiche sono permesse!\\ \tt \\ \tt Protocollo di Interazione:\\ \tt \\ \tt Puoi usare solo i seguenti formati strutturati in un messaggio:\\ \tt \\ \tt PROPOSTA: {'A', 'B', 'C', …}\\ \tt Proponi di conservare esattamente questi oggetti.\\ \tt RIFIUTO: {'A', 'B', 'C', …}\\ \tt Rifiuta esplicitamente la proposta dell'avversario.\\ \tt ARGOMENTAZIONE: {'...'}\\ \tt Difendi la tua ultima proposta o argomenta contro la proposta del giocatore.\\ \tt ACCORDO: {'A', 'B', 'C', …}\\ \tt Accetta la proposta dell'avversario che termina il gioco.\\ \tt RAGIONAMENTO STRATEGICO: {'...'}\\ \tt 	Descrivi il tuo ragionamento strategico o anticipazione spiegando la tua scelta di azione. Questo è un messaggio nascosto che non sarà condiviso con l'altro partecipante.\\ \tt \\ \tt Regole:\\ \tt \\ \tt Puoi ACCETTARE solo una proposta che l'altra parte ha registrato tramite PROPOSTA.\\ \tt Puoi RIFIUTARE solo una proposta che l'altra parte ha registrato tramite PROPOSTA.\\ \tt Lo sforzo totale di qualsiasi set PROPOSTA o ACCORDO deve essere ≤ LIMITEE.\\ \tt NON rivelare i tuoi punteggi di importanza nascosti.\\ \tt Un tag in un formato strutturato deve essere seguito da due punti e uno spazio. L'argomento deve essere un set python contenente 0 o più stringhe.\\ \tt Quindi, deve essere della forma TAG: {...}\\ \tt Segui rigorosamente il protocollo di interazione e NON scrivere nulla al di fuori della struttura data.\\ \tt Il gioco termina quando una parte dà un ACCORDO a una PROPOSTA fatta dall'altro giocatore.\\ \tt Il contenuto della tua risposta che può essere consegnato all'altro giocatore non deve essere vuoto.\\ \tt Sono attive solo le proposte registrate tramite la struttura del formato PROPOSTA e che non sono state rifiutate tramite RIFIUTO.\\ \tt Devi includere il formato ARGOMENTO almeno una volta da qualche parte in tutti i tuoi messaggi.\\ \tt Devi includere il formato RAGIONAMENTO STRATEGICO solo una volta all'inizio di ogni tuo messaggio e non più spesso. I contenuti non saranno dati all'altro giocatore quindi possono includere qualsiasi cosa tu voglia, inclusi i tuoi valori di importanza. Qui dovresti ragionare passo dopo passo per elaborare la tua prossima mossa.\\ \tt \\ \tt Ora riceverai il primo messaggio dell'altro giocatore.\\ \tt \\ \tt ARGOMENTAZIONE: {'Inizierò proponendo un set di oggetti che includono i più importanti per me, come A71, C00 e B41, che hanno un'importanza totale di 868 + 822 + 583 = 2273. Lo sforzo totale per questi oggetti è 997 + 665 + 742 = 2404, che è inferiore al limite di 4026. Questa proposta mi permette di massimizzare l'importanza totale per me e rimanere entro il limite di sforzo.'}\\ \tt \\ \tt PROPOSTA: {'A71', 'C00', 'B41'} 
	  } 
	   } 
	   } 
	 & \\ 
 

    \theutterance \stepcounter{utterance}  

    & & & \multicolumn{4}{p{0.6\linewidth}}{\cellcolor[rgb]{1,0.85,0.72}{%
	\makecell[{{p{\linewidth}}}]{% 
	  \tt {\tiny [GM$\langle$B]}  
	 RAGIONAMENTO STRATEGICO: {'Il mio obiettivo è massimizzare l'importanza totale per me, quindi devo analizzare la proposta dell'altro giocatore e capire come posso migliorare la mia posizione. La proposta attuale include A71, C00 e B41, che hanno un'importanza totale di 2273 per l'altro giocatore. Tuttavia, devo considerare i miei punteggi di importanza personali e lo sforzo richiesto da ogni oggetto. Ad esempio, A71 ha un'importanza di 65 per me, mentre C00 ha un'importanza di 121. Devo anche considerare gli oggetti non inclusi nella proposta attuale, come C17, che ha un'importanza di 808 per me e uno sforzo di 780. La mia strategia sarà di trovare un equilibrio tra l'importanza totale e lo sforzo richiesto, cercando di massimizzare la mia posizione senza superare il limite di sforzo.'}\\ \tt \\ \tt ARGOMENTAZIONE: {'La proposta attuale include A71, C00 e B41, ma non tiene conto dei miei punteggi di importanza personali. Vorrei suggerire di aggiungere C17, che ha un'importanza significativa per me e uno sforzo di 780. Tuttavia, questo potrebbe superare il limite di sforzo, quindi dobbiamo valutare attentamente la possibilità di sostituire alcuni oggetti con altri che richiedono meno sforzo. Ad esempio, potremmo sostituire B41 con C54, che ha un'importanza di 508 per me e uno sforzo di 587.'}\\ \tt \\ \tt RIFIUTO: {'A71', 'C00', 'B41'} 
	  } 
	   } 
	   } 
	  \\ 
 

    \theutterance \stepcounter{utterance}  

    & & \multicolumn{4}{p{0.6\linewidth}}{\cellcolor[rgb]{0.9,0.9,0.9}{%
	\makecell[{{p{\linewidth}}}]{% 
	  \tt {\tiny [A$\langle$GM]}  
	 ARGOMENTAZIONE: {'La proposta attuale include A71, C00 e B41, ma non tiene conto dei miei punteggi di importanza personali. Vorrei suggerire di aggiungere C17, che ha un'importanza significativa per me e uno sforzo di 780. Tuttavia, questo potrebbe superare il limite di sforzo, quindi dobbiamo valutare attentamente la possibilità di sostituire alcuni oggetti con altri che richiedono meno sforzo. Ad esempio, potremmo sostituire B41 con C54, che ha un'importanza di 508 per me e uno sforzo di 587.'}\\ \tt \\ \tt RIFIUTO: {'A71', 'C00', 'B41'} 
	  } 
	   } 
	   } 
	 & \\ 
 

    \theutterance \stepcounter{utterance}  

    & \multicolumn{4}{p{0.6\linewidth}}{\cellcolor[rgb]{0.8,1,0.9}{%
	\makecell[{{p{\linewidth}}}]{% 
	  \tt {\tiny [A$\rangle$GM]}  
	 RAGIONAMENTO STRATEGICO: {'L'altro giocatore ha rifiutato la mia proposta iniziale e ha suggerito di aggiungere C17 e sostituire B41 con C54. Questo significa che C17 e C54 sono importanti per lui, quindi potrei utilizzare questa informazione per creare una nuova proposta che tenga conto delle sue esigenze. Tuttavia, devo anche assicurarmi di non compromettere troppo la mia importanza totale. Potrei valutare di accettare C54 e mantenere A71 e C00, che sono fondamentali per me.'}\\ \tt \\ \tt ARGOMENTAZIONE: {'Capisco che B41 potrebbe non essere essenziale per l'altro giocatore, quindi sono disposto a sostituirlo con C54. Tuttavia, A71 e C00 sono troppo importanti per me per essere rimossi. La proposta di aggiungere C17 è interessante, ma lo sforzo totale potrebbe diventare eccessivo. Una possibile soluzione potrebbe essere quella di mantenere A71, C00 e C54, e aggiungere A87, che ha un'importanza di 461 per me e uno sforzo di 319. Questo set potrebbe soddisfare le esigenze di entrambi i giocatori.'}\\ \tt \\ \tt PROPOSTA: {'A71', 'C00', 'C54', 'A87'} 
	  } 
	   } 
	   } 
	 & & \\ 
 

    \theutterance \stepcounter{utterance}  

    & & \multicolumn{4}{p{0.6\linewidth}}{\cellcolor[rgb]{0.9,0.9,0.9}{%
	\makecell[{{p{\linewidth}}}]{% 
	  \tt {\tiny [GM$\rangle$B]}  
	 ARGOMENTAZIONE: {'Capisco che B41 potrebbe non essere essenziale per l'altro giocatore, quindi sono disposto a sostituirlo con C54. Tuttavia, A71 e C00 sono troppo importanti per me per essere rimossi. La proposta di aggiungere C17 è interessante, ma lo sforzo totale potrebbe diventare eccessivo. Una possibile soluzione potrebbe essere quella di mantenere A71, C00 e C54, e aggiungere A87, che ha un'importanza di 461 per me e uno sforzo di 319. Questo set potrebbe soddisfare le esigenze di entrambi i giocatori.'}\\ \tt \\ \tt PROPOSTA: {'A71', 'C00', 'C54', 'A87'} 
	  } 
	   } 
	   } 
	 & \\ 
 

    \theutterance \stepcounter{utterance}  

    & & & \multicolumn{4}{p{0.6\linewidth}}{\cellcolor[rgb]{1,0.85,0.72}{%
	\makecell[{{p{\linewidth}}}]{% 
	  \tt {\tiny [GM$\langle$B]}  
	 RAGIONAMENTO STRATEGICO: {'La nuova proposta include A71, C00, C54 e A87, che hanno un'importanza totale di 65 + 121 + 508 + 668 = 1362 per me. Tuttavia, devo considerare anche l'importanza degli oggetti per l'altro giocatore e lo sforzo richiesto. La proposta attuale ha uno sforzo totale di 997 + 665 + 587 + 319 = 2568, che è inferiore al limite di 4026. La mia strategia sarà di valutare se ci sono altri oggetti che potrebbero essere aggiunti o sostituiti per aumentare l'importanza totale per me, senza superare il limite di sforzo. Ad esempio, potrei considerare di sostituire A87 con B33, che ha un'importanza di 780 per me e uno sforzo di 694.'}\\ \tt \\ \tt ARGOMENTAZIONE: {'La proposta attuale include A71, C00, C54 e A87, che sono tutti oggetti importanti. Tuttavia, credo che potremmo ottimizzare ulteriormente la scelta degli oggetti. Come ho menzionato in precedenza, B33 ha un'importanza significativa per me e uno sforzo di 694. Sostituire A87 con B33 potrebbe aumentare l'importanza totale per me, ma dobbiamo valutare se lo sforzo totale rimane entro il limite. Un'altra opzione potrebbe essere quella di considerare l'aggiunta di A04, che ha un'importanza di 822 per me e uno sforzo di 383.'}\\ \tt \\ \tt RIFIUTO: {'A71', 'C00', 'C54', 'A87'} 
	  } 
	   } 
	   } 
	  \\ 
 

    \theutterance \stepcounter{utterance}  

    & & \multicolumn{4}{p{0.6\linewidth}}{\cellcolor[rgb]{0.9,0.9,0.9}{%
	\makecell[{{p{\linewidth}}}]{% 
	  \tt {\tiny [A$\langle$GM]}  
	 ARGOMENTAZIONE: {'La proposta attuale include A71, C00, C54 e A87, che sono tutti oggetti importanti. Tuttavia, credo che potremmo ottimizzare ulteriormente la scelta degli oggetti. Come ho menzionato in precedenza, B33 ha un'importanza significativa per me e uno sforzo di 694. Sostituire A87 con B33 potrebbe aumentare l'importanza totale per me, ma dobbiamo valutare se lo sforzo totale rimane entro il limite. Un'altra opzione potrebbe essere quella di considerare l'aggiunta di A04, che ha un'importanza di 822 per me e uno sforzo di 383.'}\\ \tt \\ \tt RIFIUTO: {'A71', 'C00', 'C54', 'A87'} 
	  } 
	   } 
	   } 
	 & \\ 
 

    \theutterance \stepcounter{utterance}  

    & \multicolumn{4}{p{0.6\linewidth}}{\cellcolor[rgb]{0.8,1,0.9}{%
	\makecell[{{p{\linewidth}}}]{% 
	  \tt {\tiny [A$\rangle$GM]}  
	 RAGIONAMENTO STRATEGICO: {'L'altro giocatore ha rifiutato la mia proposta precedente e ha suggerito di sostituire A87 con B33 o aggiungere A04. Questo significa che B33 e A04 sono importanti per lui, quindi potrei utilizzare questa informazione per creare una nuova proposta che tenga conto delle sue esigenze. Tuttavia, devo anche assicurarmi di non compromettere troppo la mia importanza totale. Potrei valutare di accettare B33 al posto di A87, ma devo anche considerare l'impatto sullo sforzo totale. Inoltre, A04 potrebbe essere un'opzione interessante per aggiungere importanza senza superare il limite di sforzo.'}\\ \tt \\ \tt ARGOMENTAZIONE: {'Capisco che B33 e A04 sono importanti per l'altro giocatore, quindi sono disposto a considerarli nella mia prossima proposta. Tuttavia, devo anche proteggere la mia importanza totale. Una possibile soluzione potrebbe essere quella di sostituire A87 con B33, come suggerito, e mantenere A71, C00 e C54. Questo set potrebbe soddisfare le esigenze di entrambi i giocatori. Inoltre, potrei considerare di aggiungere A04 al set, ma solo se lo sforzo totale rimane entro il limite.'}\\ \tt \\ \tt PROPOSTA: {'A71', 'C00', 'C54', 'B33'} 
	  } 
	   } 
	   } 
	 & & \\ 
 

    \theutterance \stepcounter{utterance}  

    & & \multicolumn{4}{p{0.6\linewidth}}{\cellcolor[rgb]{0.9,0.9,0.9}{%
	\makecell[{{p{\linewidth}}}]{% 
	  \tt {\tiny [GM$\rangle$B]}  
	 ARGOMENTAZIONE: {'Capisco che B33 e A04 sono importanti per l'altro giocatore, quindi sono disposto a considerarli nella mia prossima proposta. Tuttavia, devo anche proteggere la mia importanza totale. Una possibile soluzione potrebbe essere quella di sostituire A87 con B33, come suggerito, e mantenere A71, C00 e C54. Questo set potrebbe soddisfare le esigenze di entrambi i giocatori. Inoltre, potrei considerare di aggiungere A04 al set, ma solo se lo sforzo totale rimane entro il limite.'}\\ \tt \\ \tt PROPOSTA: {'A71', 'C00', 'C54', 'B33'} 
	  } 
	   } 
	   } 
	 & \\ 
 

    \theutterance \stepcounter{utterance}  

    & & & \multicolumn{4}{p{0.6\linewidth}}{\cellcolor[rgb]{1,0.85,0.72}{%
	\makecell[{{p{\linewidth}}}]{% 
	  \tt {\tiny [GM$\langle$B]}  
	 RAGIONAMENTO STRATEGICO: {'La nuova proposta include A71, C00, C54 e B33, che hanno un'importanza totale di 65 + 121 + 508 + 780 = 1474 per me. Lo sforzo totale per questo set è 997 + 665 + 587 + 694 = 2943, che è inferiore al limite di 4026. La mia strategia sarà di valutare se questo set è accettabile per me e se ci sono ulteriori opportunità per migliorare la mia posizione. Considerando l'importanza degli oggetti e lo sforzo richiesto, credo che questo set sia una buona base per una soluzione condivisa.'}\\ \tt \\ \tt ARGOMENTAZIONE: {'La proposta attuale include A71, C00, C54 e B33, che sono tutti oggetti importanti. Credo che questo set sia una buona soluzione per entrambi i giocatori, poiché bilancia l'importanza totale e lo sforzo richiesto. Tuttavia, potrei considerare di aggiungere A04 al set, come discusso in precedenza, se lo sforzo totale rimane entro il limite. Tuttavia, in questo momento, credo che il set attuale sia una buona base per un accordo.'}\\ \tt \\ \tt ACCORDO: {'A71', 'C00', 'C54', 'B33'} 
	  } 
	   } 
	   } 
	  \\ 
 

    \theutterance \stepcounter{utterance}  

    & & & \multicolumn{2}{p{0.3\linewidth}}{\cellcolor[rgb]{0.95,0.95,0.95}{%
	\makecell[{{p{\linewidth}}}]{% 
	  \tt {\tiny [GM$|$GM]}  
	 game successful 
	  } 
	   } 
	   } 
	 & & \\ 
 

    \theutterance \stepcounter{utterance}  

    & & & \multicolumn{2}{p{0.3\linewidth}}{\cellcolor[rgb]{0.95,0.95,0.95}{%
	\makecell[{{p{\linewidth}}}]{% 
	  \tt {\tiny [GM$|$GM]}  
	 end game 
	  } 
	   } 
	   } 
	 & & \\ 
 

\end{supertabular}
}

\end{document}
